\documentclass[12pt]{article}
\usepackage[T1]{fontenc}
\usepackage[utf8]{inputenc}
\usepackage{lmodern}
\usepackage{textcomp}
\usepackage{enumitem}
\usepackage{geometry}
\geometry{margin=1in}
\begin{document}
\begin{center}
Answer Key - Astronomy - Graduate Level Assessment 2 \
\small (Answers are ordered exactly as in the question paper.)
\end{center}

\section*{Multiple Choice}
\begin{enumerate}[label=\arabic*.]
\item \textbf{Answer:} B
\\ \textit{Options:} A) exactly 3 times longer; B) about 5.2 times longer; C) Not enough information. It will depend on the inclination of the new orbit; D) The length of the year wouldn't change because the Earth's mass stays the same.
\\ \textit{Note:} The correct answer is based on Kepler's Third Law of planetary motion, which states that the square of the orbital period of a planet is directly proportional to the cube of the semi-major axis of its orbit, meaning if the Earth were three times further from the Sun, its orbital period would increase by a factor of approximately 5.2.
\item \textbf{Answer:} A
\\ \textit{Options:} A) 10000 times more; B) 100 times more; C) 1000 times more; D) 10 times more
\\ \textit{Note:} The telescope can gather light more effectively due to its larger aperture, which collects light proportional to the square of the diameter; thus, the ratio of the light-gathering power is (500 mm / 5 mm)² = 10000.
\item \textbf{Answer:} A
\\ \textit{Options:} A) The Moon split from the Earth due to tidal forces.; B) The Moon was captured into Earth orbit.; C) The Earth and Moon co-accreted in the solar nebula.; D) Earth was rotating so rapidly that the Moon split from it.
\\ \textit{Note:} The hypothesis that the Moon split from the Earth due to tidal forces is not widely supported in the scientific community, as it contradicts current understanding of lunar formation, which favors theories like the giant impact hypothesis, fission, and co-formation.
\item \textbf{Answer:} D
\\ \textit{Options:} A) Perseus; B) Cygnus; C) Scorpius; D) Leo
\\ \textit{Note:} Leo is a zodiac constellation that is situated far from the plane of the Milky Way, unlike other constellations such as Scorpius or Sagittarius, which are prominently located along the galactic band.
\item \textbf{Answer:} D
\\ \textit{Options:} A) The Earth must lie completely within the Moon's penumbra.; B) The Moon's penumbra must touch the area where you are located.; C) The Earth must be near aphelion in its orbit of the Sun.; D) The Moon's umbra must touch the area where you are located.
\\ \textit{Note:} To observe a total solar eclipse, the Moon's umbra must reach the Earth at your location, as only within this shadow can the Sun be completely obscured by the Moon, allowing for the totality phase of the eclipse.
\end{enumerate}

\section*{True / False}
\begin{enumerate}[label=\arabic*.]
\setcounter{enumi}{5}
\item \textbf{Answer:} True
\\ \textit{Options:} A) True; B) False
\\ \textit{Note:} The statement is true because the craters at Mercury's poles are located in such a way that they never receive direct sunlight, creating permanently shadowed regions that are cold enough to allow water-ice to persist.
\item \textbf{Answer:} True
\\ \textit{Options:} A) True; B) False
\\ \textit{Note:} The Andromeda Galaxy is located at a distance of about 2.5 million light years from Earth, as confirmed by various astronomical measurements and observations, making the statement true.
\item \textbf{Answer:} False
\\ \textit{Options:} A) True; B) False
\\ \textit{Note:} The statement is false because Mars has a weaker gravitational pull than Earth, meaning that the truck's weight would actually be less on Mars, making it easier to accelerate.
\item \textbf{Answer:} False
\\ \textit{Options:} A) True; B) False
\\ \textit{Note:} The resolution of a telescope is primarily determined by its aperture size and optical quality, rather than its magnification capability, as resolution refers to the telescope's ability to distinguish between closely spaced objects.
\item \textbf{Answer:} False
\\ \textit{Options:} A) True; B) False
\\ \textit{Note:} The statement is false because the Moon's surface is more heavily cratered than the Earth's due to the lack of significant geological activity and erosion processes that would otherwise erase impact craters on Earth.
\end{enumerate}

\section*{Long Form Response}
\begin{enumerate}[label=\arabic*.]
\setcounter{enumi}{10}
\item \textbf{Answer:} The compositional differences between the inner and outer planets in our solar system can be primarily attributed to the varying temperatures present during the condensation of materials in the protoplanetary nebula. In the hotter inner regions, only metals and silicate rocks could condense, leading to the formation of terrestrial planets with solid surfaces, such as Mercury, Venus, Earth, and Mars. Conversely, in the cooler outer regions of the nebula, hydrogen compounds, such as water, ammonia, and methane, were able to condense, allowing for the formation of gas giants like Jupiter and Saturn. This temperature gradient thus played a crucial role in determining the materials available for planet formation, ultimately influencing the distinct compositions of the inner and outer planets. As a result, the inner planets are rocky and metallic, while the outer planets are predominantly gaseous with significant amounts of ices.
\\ \textit{Note:} correct because it accurately describes how temperature gradients in the protoplanetary nebula influenced the condensation of materials, resulting in the formation of rocky terrestrial planets in the inner solar system and gas giant planets with volatile compounds in the cooler outer regions.
\item \textbf{Answer:} Bolometric luminosity is a critical concept in astronomy that refers to the total energy output of a star or astronomical object integrated over all wavelengths of electromagnetic radiation, from radio waves to gamma rays. This measure is significant because it provides a comprehensive understanding of an object's energy production, allowing astronomers to assess its intrinsic brightness regardless of distance or intervening materials that may obscure certain wavelengths. In the context of stellar observations, bolometric luminosity is calculated by integrating the spectral energy distribution (SED) across all wavelengths, often requiring corrections for the effects of interstellar extinction and the object's distance. By evaluating bolometric luminosity, astronomers can derive important physical parameters of stars, such as temperature and radius, and make comparisons across different stellar populations.
\\ \textit{Note:} The answer correctly identifies bolometric luminosity as the total energy output of a star across all electromagnetic wavelengths, emphasizing its importance for understanding intrinsic brightness and energy production while addressing the impact of distance and obscuration on stellar observations.
\end{enumerate}

\vfill
\noindent\textit{End of Answer Key}
\end{document}